% Archive record 14, line 1.
\hypertarget{part-2-14-0001-ii-b-concentration-inequalities-and-uniform-convergence}{%
\section{II-B --- Concentration Inequalities and Uniform Convergence}\label{part-2-14-0001-ii-b-concentration-inequalities-and-uniform-convergence}}

上一阶段 II-A，我们得到了一个重要结论：
\[R(\hat h_n)-R(h^*)
\leq
2\sup_{h\in\mathcal H}
|\widehat R_n(h)-R(h)|.\]

也就是说，只要能够控制整个 Hypothesis Space 中最大的评估误差，就能够控制 Empirical Risk Minimization 所选择模型的实际性能损失。

但我们还不知道这个误差究竟有多大。

本阶段的目标，就是建立 \textbf{Sample Size、Search Space 和 Generalization Error 之间的定量关系}。

\hypertarget{part-2-14-0001-knowledge-roadmap}{%
\subsection{1. Knowledge Roadmap}\label{part-2-14-0001-knowledge-roadmap}}

\textbf{Law of Large Numbers}

当样本量趋于无穷时，误差是否消失？

↓ \textbf{Concentration Inequalities}

有限样本下，误差超过阈值的概率有多大？

↓ \textbf{Hoeffding's Inequality}

控制单个固定模型的评估误差

↓ \textbf{Union Bound}

同时控制多个候选模型的评估误差

↓ \textbf{Uniform Generalization Bound}

控制整个假设空间，进而保证模型选择的可靠性

前两个部分与你已经学习过的 Probability 内容高度重合，我们快速通过，重点关注后三者如何应用于 Statistical Learning。

\hypertarget{part-2-14-0001-from-lln-to-concentration-inequalities}{%
\subsection{2. From LLN to Concentration Inequalities}\label{part-2-14-0001-from-lln-to-concentration-inequalities}}

\hypertarget{part-2-14-0001-lln-ux7684ux5c40ux9650ux6027}{%
\subsubsection{2.1 LLN 的局限性}\label{part-2-14-0001-lln-ux7684ux5c40ux9650ux6027}}

沿用 II-A 中的记号：
\[R(h)=\mathbb E[\ell(h(X),Y)],\]

\[\widehat R_n(h)
=
\frac1n\sum_{i=1}^{n}\ell(h(X_i),Y_i).\]

对于一个事先固定的模型 \(h\)，在 IID Sampling 和损失可积的条件下，根据 SLLN：
\[\widehat R_n(h)
\xrightarrow{\mathrm{a.s.}}
R(h).\]

然而，LLN 是一个 \textbf{Asymptotic Result（渐近结论）}。

它告诉我们当 \(n\to\infty\) 时会发生什么，但并不直接告诉我们：

\begin{itemize}
\tightlist
\item
  当 \(n=500\) 时，评估结果有多可靠？
\item
  我们至少需要多少样本才能控制评估误差？
\item
  当 Agent 比较大量候选模型时，需要增加多少评估数据？
\end{itemize}

因此，我们希望得到如下形式的 \textbf{Finite-Sample Guarantee}：
\[\boxed{
\mathbb P
\left(
|\widehat R_n(h)-R(h)|>\epsilon
\right)
\leq B(n,\epsilon)
}\]

其中 \(B(n,\epsilon)\) 是一个可以明确计算的概率上界。

这就是 Concentration Inequalities 的作用。

\hypertarget{part-2-14-0001-ux5e38ux89c1-concentration-inequalities}{%
\subsubsection{2.2 常见 Concentration Inequalities}\label{part-2-14-0001-ux5e38ux89c1-concentration-inequalities}}

设随机变量 \(X_1,\ldots,X_n\) IID，均值为 \(\mu\)，方差为 \(\sigma^2\)，并定义：
\[\bar X_n=\frac1n\sum_{i=1}^{n}X_i.\]

\begin{readertable}
\begin{longtable}[]{@{}
  >{\raggedright\arraybackslash}p{(\columnwidth - 4\tabcolsep) * \real{0.33}}
  >{\raggedright\arraybackslash}p{(\columnwidth - 4\tabcolsep) * \real{0.33}}
  >{\raggedright\arraybackslash}p{(\columnwidth - 4\tabcolsep) * \real{0.33}}@{}}
\toprule
Inequality & 主要条件 & Tail Probability Bound \\
\midrule
\endhead
Markov & 非负随机变量 \(X\) & \(\mathbb P(X\geq a)\leq\frac{\mathbb E[X]}a\) \\
Chebyshev & 有限方差 & \(\mathbb P(|\bar X_n-\mu|\geq\epsilon)\leq\frac{\sigma^2}{n\epsilon^2}\) \\
Hoeffding & 独立、有界变量 \(X_i\in[0,1]\) & \(2e^{-2n\epsilon^2}\) \\
Bernstein & 独立、有界，利用方差信息 & 指数型上界，考虑实际方差 \\
\bottomrule
\end{longtable}
\end{readertable}

最值得关注的是 Chebyshev 与 Hoeffding 在收敛速度上的区别。

对于固定的 \(\epsilon>0\)：
\[B_{\mathrm{Chebyshev}}(n,\epsilon)
=
O(n^{-1}),\]

而：
\[B_{\mathrm{Hoeffding}}(n,\epsilon)
=
O(e^{-cn}),
\qquad c=2\epsilon^2.\]

\textbf{Chebyshev vs. Hoeffding}

假设 IID 随机变量取值于 {[}0,1{]}，Chebyshev 使用最大可能方差 1/4。

Error tolerance \(\epsilon\)\\
\textbf{0.10}

0.1（范围 0.05--0.2）

\bookgraphic{0.498}{part-2-14-0001-002-01}{}

Chebyshev

Hoeffding

纵轴采用 Log Scale；概率上界截断至 1，低于 10⁻⁶ 的部分不显示。

eps = 0.2

\bookgraphic{0.498}{part-2-14-0001-002-state-01}{}

\textbf{核心观点：} Hoeffding 利用更强的分布条件，换取了更强的有限样本误差控制。

\hypertarget{part-2-14-0001-hoeffdings-inequality}{%
\subsection{3. Hoeffding's Inequality}\label{part-2-14-0001-hoeffdings-inequality}}

\hypertarget{part-2-14-0001-theorem}{%
\subsubsection{3.1 Theorem}\label{part-2-14-0001-theorem}}

设 \(X_1,\ldots,X_n\) 为相互独立的随机变量，且：
\[a_i\leq X_i\leq b_i
\qquad\text{a.s.}\]

令：
\[S_n=\sum_{i=1}^{n}X_i.\]

则对于任意 \(t>0\)：
\[\boxed{
\mathbb P
\left(
|S_n-\mathbb E[S_n]|\geq t
\right)
\leq
2\exp
\left(
-\frac{2t^2}
{\sum_{i=1}^{n}(b_i-a_i)^2}
\right)
}\]

特别地，如果：
\[X_i\in[0,1],\]

则：
\[\boxed{
\mathbb P
\left(
|\bar X_n-\mathbb E[\bar X_n]|
\geq\epsilon
\right)
\leq
2e^{-2n\epsilon^2}
}\]

注意，Hoeffding 不要求随机变量服从 Normal Distribution，也不要求它们具有完全相同的分布。

\hypertarget{part-2-14-0001-proof-index-mathematical-principle}{%
\subsubsection{3.2 Proof Index --- Mathematical Principle}\label{part-2-14-0001-proof-index-mathematical-principle}}

完整推导目前不是重点。它主要依赖于你已经接触过的 MGF、Markov's Inequality 和 Independence。

\textbf{Hoeffding's Lemma} \(\mathbb E[e^{\lambda(X-\mathbb E[X])}]\leq e^{\lambda^2(b-a)^2/8}\) ↓ \textbf{Independence}

将随机变量之和的 MGF 分解为乘积

↓ \textbf{Exponential Markov Inequality} \(\mathbb P(S_n-\mathbb E[S_n]\geq t)\leq e^{-\lambda t}\mathbb E[e^{\lambda(S_n-\mathbb E[S_n])}]\) ↓ \textbf{Optimize over \(\lambda>0\)}

得到最紧的指数型上界，再合并两侧尾概率

现阶段掌握定理本身以及其成立条件即可。

\hypertarget{part-2-14-0001-applying-hoeffding-to-statistical-learning}{%
\subsection{4. Applying Hoeffding to Statistical Learning}\label{part-2-14-0001-applying-hoeffding-to-statistical-learning}}

\hypertarget{part-2-14-0001-single-model-generalization-bound}{%
\subsubsection{4.1 Single-model Generalization Bound}\label{part-2-14-0001-single-model-generalization-bound}}

现在将随机变量替换为模型的 Loss：
\[L_i(h)=\ell(h(X_i),Y_i).\]

假设：

\begin{enumerate}
\def\labelenumi{\arabic{enumi}.}
\tightlist
\item
  \((X_i,Y_i)\) IID；
\item
  \(h\) 事先固定且独立于评估数据；
\item
  Loss Function 满足：
\end{enumerate}

\[0\leq\ell(h(X_i),Y_i)\leq1.\]

那么：
\[\mathbb E[L_i(h)]=R(h),\]

\[\frac1n\sum_{i=1}^{n}L_i(h)
=
\widehat R_n(h).\]

直接应用 Hoeffding：
\[\boxed{
\mathbb P
\left(
|\widehat R_n(h)-R(h)|>\epsilon
\right)
\leq
2e^{-2n\epsilon^2}
}\]

这就是 \textbf{Single-model Generalization Bound}。

\hypertarget{part-2-14-0001-confidence-form}{%
\subsubsection{4.2 Confidence Form}\label{part-2-14-0001-confidence-form}}

引入 Failure Probability：
\[\alpha\in(0,1).\]

这里使用 \(\alpha\)，以免与 II-A 中表示 Optimization Error 的 \(\delta\) 混淆。

令：
\[2e^{-2n\epsilon^2}=\alpha.\]

解得：
\[\epsilon
=
\sqrt{\frac{\log(2/\alpha)}{2n}}.\]

因此，以至少 \(1-\alpha\) 的概率：
\[\boxed{
|\widehat R_n(h)-R(h)|
\leq
\sqrt{\frac{\log(2/\alpha)}{2n}}
}\]

这已经建立了第一个重要的定量关系：
\[\text{Generalization Error Bound}
=
O\left(\frac1{\sqrt n}\right).\]

即使我们完全不知道 Loss 的具体分布，也能够在有界性和独立性条件下给出这一保证。

但目前仍然只能控制一个固定模型。

\hypertarget{part-2-14-0001-union-bound-and-uniform-generalization}{%
\subsection{5. Union Bound and Uniform Generalization}\label{part-2-14-0001-union-bound-and-uniform-generalization}}

\textbf{这是本阶段最重要的内容。}

\hypertarget{part-2-14-0001-research-motivation}{%
\subsubsection{5.1 Research Motivation}\label{part-2-14-0001-research-motivation}}

假设 RSI Agent 并不是只评价一个模型，而是从以下候选集合中选择模型：
\[\mathcal H=\{h_1,h_2,\ldots,h_M\}.\]

其中：
\[M=|\mathcal H|.\]

即使每个模型的评估误差超过 \(\epsilon\) 的概率都很小，只要候选模型数量足够多，仍然可能出现某一个模型具有很大的评估误差。

而 Agent 恰恰可能选择到这个模型。

因此，我们需要从控制单个模型：
\[|\widehat R_n(h)-R(h)|\]

转向同时控制所有模型：
\[\sup_{h\in\mathcal H}
|\widehat R_n(h)-R(h)|.\]

\hypertarget{part-2-14-0001-theorem-finite-hypothesis-space-generalization-bound}{%
\subsubsection{5.2 Theorem --- Finite Hypothesis Space Generalization Bound}\label{part-2-14-0001-theorem-finite-hypothesis-space-generalization-bound}}

假设：

\begin{itemize}
\tightlist
\item
  \(\mathcal H\) 是事先确定的有限假设空间；
\item
  评估样本 IID；
\item
  所有模型的 Loss 均取值于 \([0,1]\)。
\end{itemize}

那么：
\[\boxed{
\mathbb P
\left(
\sup_{h\in\mathcal H}
|\widehat R_n(h)-R(h)|>\epsilon
\right)
\leq
2Me^{-2n\epsilon^2}
}\]

这就是我们一直在追求的 \textbf{Uniform Generalization Bound}。

\hypertarget{part-2-14-0001-proof}{%
\subsubsection{5.3 Proof}\label{part-2-14-0001-proof}}

这个证明非常简洁，也值得理解。

定义事件：
\[E_j=
\left\{
|\widehat R_n(h_j)-R(h_j)|>\epsilon
\right\}.\]

整个假设空间中至少有一个模型的误差超过 \(\epsilon\)，等价于：
\[\bigcup_{j=1}^{M}E_j.\]

因此，根据 Union Bound：
\[\begin{aligned}
&\mathbb P
\left(
\sup_{h\in\mathcal H}
|\widehat R_n(h)-R(h)|>\epsilon
\right)\\
&=
\mathbb P\left(\bigcup_{j=1}^{M}E_j\right)\\
&\leq
\sum_{j=1}^{M}\mathbb P(E_j)\\
&\leq
\sum_{j=1}^{M}2e^{-2n\epsilon^2}\\
&=
2Me^{-2n\epsilon^2}.
\end{aligned}\]

\[\hfill\blacksquare\]

有一个非常重要的细节：

\textbf{Union Bound 不要求不同模型的评估误差相互独立。}

因此，即使 RSI 产生的多个 Agent 版本具有高度相似的结构，只要它们属于满足上述条件的固定有限假设空间，该不等式仍然成立。

\hypertarget{part-2-14-0001-confidence-form-1}{%
\subsubsection{5.4 Confidence Form}\label{part-2-14-0001-confidence-form-1}}

与前面采用相同的方法，令：
\[2Me^{-2n\epsilon^2}=\alpha.\]

解得：
\[\boxed{
\epsilon
=
\sqrt{
\frac{\log(2M/\alpha)}{2n}
}
}\]

因此，以至少 \(1-\alpha\) 的概率：
\[\boxed{
\forall h\in\mathcal H,\quad
|\widehat R_n(h)-R(h)|
\leq
\sqrt{
\frac{\log(2M/\alpha)}{2n}
}
}\]

这个式子非常重要，因为它直接建立了四个变量之间的联系。

\textbf{Uniform Generalization Bound Explorer}

\[\epsilon=\sqrt{\frac{\log(2M/\alpha)}{2n}}\]

Evaluation Sample Size \(n\)\\
\textbf{1,000}

3（范围 2--5）

Hypothesis Space Size \(M\)\\
\textbf{1,000}

3（范围 0--6） Confidence Level \(1-\alpha\) 90\% / 95\% / 99\%

Uniform Error Bound

\textbf{0.0728}

ERM Excess Risk Bound

\textbf{0.1456}

\bookgraphic{0.588}{part-2-14-0001-004-01}{该工具计算的是理论概率上界，而不是对真实 Generalization Gap 的预测。}

logN = 5

\bookgraphic{0.588}{part-2-14-0001-004-state-01}{}

我们可以直接观察到：

\begin{readertable}
\begin{longtable}[]{@{}
  >{\raggedright\arraybackslash}p{(\columnwidth - 4\tabcolsep) * \real{0.18}}
  >{\raggedright\arraybackslash}p{(\columnwidth - 4\tabcolsep) * \real{0.46}}
  >{\raggedright\arraybackslash}p{(\columnwidth - 4\tabcolsep) * \real{0.32}}@{}}
\toprule
Variable & Change & Impact on Error Bound \\
\midrule
\endhead
\(n\) & 样本量增加 & 下降 \\
\(M\) & 假设空间扩大 & 上升 \\
\(\alpha\) & 要求更高置信水平，即 \(\alpha\) 减小 & 上升 \\
\(\epsilon\) & 要求更小的误差 & 需要更多数据 \\
\bottomrule
\end{longtable}
\end{readertable}

尤其注意：
\[\boxed{
\text{Uniform Generalization Error Bound}
=
O\left(
\sqrt{\frac{\log M+\log(1/\alpha)}n}
\right)
}\]

这与 Part I 中通过 Gaussian Maximum 得到的 Selection Bias 尺度：
\[O\left(\sqrt{\frac{\log M}{n}}\right)\]

具有非常相似的数学形式。

但两者不是同一个结论。前者是有界损失条件下的高概率上界，后者是特定噪声模型下的期望偏差尺度。

\hypertarget{part-2-14-0001-sample-complexity}{%
\subsection{6. Sample Complexity}\label{part-2-14-0001-sample-complexity}}

我们也可以从反方向理解上述不等式。

假设我们要求：
\[\sup_{h\in\mathcal H}
|\widehat R_n(h)-R(h)|
\leq\epsilon\]

以至少 \(1-\alpha\) 的概率成立。

那么，一个充分的样本量条件是：
\[\boxed{
n\geq
\frac{\log(2M/\alpha)}{2\epsilon^2}
}\]

这就是有限假设空间下的一种 \textbf{Sample Complexity Bound}。

\textbf{Required Evaluation Sample Size}

固定 Confidence Level = 95\%

Error Tolerance \(\epsilon\)\\
\textbf{0.05}

0.05（范围 0.01--0.2）

Hypothesis Space Size \(M\)\\
\textbf{1,000}

3（范围 0--6）

该理论保证所需的充分样本量

\textbf{2,120}

将误差容忍度减半，充分样本量约增加至原来的四倍。

最值得记住的是两个 Scaling Laws：
\[n=O(\log M)\]

以及：
\[n=O(\epsilon^{-2}).\]

这里分别固定其他变量。

这意味着，在该理论框架下，扩大有限搜索空间所增加的样本需求是对数级别的，而提高评估精度所增加的样本需求是平方反比级别的。

不过，这是一个充分条件，并不是实际任务所需最小样本量的精确计算。

\hypertarget{part-2-14-0001-combining-uniform-convergence-with-excess-risk}{%
\subsection{7. Combining Uniform Convergence with Excess Risk}\label{part-2-14-0001-combining-uniform-convergence-with-excess-risk}}

现在我们将 II-A 与 II-B 的数学结论连接起来。

II-A 中已经证明：
\[R(\hat h_n)-R(h^*)
\leq
2\sup_{h\in\mathcal H}
|\widehat R_n(h)-R(h)|.\]

II-B 中又证明，以至少 \(1-\alpha\) 的概率：
\[\sup_{h\in\mathcal H}
|\widehat R_n(h)-R(h)|
\leq
\sqrt{\frac{\log(2M/\alpha)}{2n}}.\]

将两个结果结合：
\[\boxed{
R(\hat h_n)-R(h^*)
\leq
\sqrt{
\frac{2\log(2M/\alpha)}{n}
}
}\]

这就是有限假设空间下 ERM 的一种 \textbf{Excess Risk Guarantee}。

对于 Approximate ERM，进一步有：
\[\boxed{
R(\tilde h_n)-R(h^*)
\leq
\underbrace{
\sqrt{\frac{2\log(2M/\alpha)}n}
}_{\text{Statistical Error Bound}}
+
\underbrace{\delta}_{\text{Optimization Error Bound}}
}\]

这是目前为止我们得到的、与 RSI 研究联系最紧密的数学结论之一。

\hypertarget{part-2-14-0001-optimizationgeneralization-trade-off}{%
\subsubsection{7.1 Optimization--Generalization Trade-off}\label{part-2-14-0001-optimizationgeneralization-trade-off}}

假设 Agent 的改进能力不断增强：

\begin{itemize}
\tightlist
\item
  它能够搜索更多候选方法；
\item
  它更容易找到 Empirical Risk 较低的模型；
\item
  它的 Optimization Error 可能不断下降；
\item
  但如果有效搜索空间扩大，Statistical Error Bound 可能上升。
\end{itemize}

\textbf{Optimization--Generalization Trade-off}

\bookgraphic{0.498}{part-2-14-0001-006-01}{}

Optimization error

Statistical error bound

Combined bound

研究假设示意：假设搜索能力增强的同时，有效搜索空间也在扩大。该曲线不是上述定理必然推出的结果。

这里我们已经获得了一个解释 Better, Then Worse 的潜在数学框架。

但必须强调：\textbf{上界变差不等于真实性能变差。} 我们还不能据此证明 RSI 必然产生 Performance Decay。

\hypertarget{part-2-14-0001-limitations-ux4e3aux4ec0ux4e48ux8fd9ux4e9bux5b9aux7406ux4e0dux80fdux76f4ux63a5ux5957ux7528ux5230-rsi}{%
\subsection{8. Limitations --- 为什么这些定理不能直接套用到 RSI？}\label{part-2-14-0001-limitations-ux4e3aux4ec0ux4e48ux8fd9ux4e9bux5b9aux7406ux4e0dux80fdux76f4ux63a5ux5957ux7528ux5230-rsi}}

这一节在科研意义上甚至比定理本身更重要。

\hypertarget{part-2-14-0001-fixed-hypothesis-space-vs.-adaptive-hypothesis-space}{%
\subsubsection{8.1 Fixed Hypothesis Space vs. Adaptive Hypothesis Space}\label{part-2-14-0001-fixed-hypothesis-space-vs.-adaptive-hypothesis-space}}

在前面的证明中，我们假设：
\[\mathcal H=\{h_1,\ldots,h_M\}\]

是一个事先确定的假设空间。

但是，在 RSI 中，候选版本可能根据之前的评估结果不断生成：
\[A_{k+1}
=
\mathcal I_{A_k}(A_k,F_k,H_k).\]

因此，实际访问的候选集合可能依赖于评估数据 \(D_n\)：
\[\mathcal H(D_n)
=
\{A_0,A_1(D_n),\ldots,A_K(D_n)\}.\]

\textbf{Fixed Search Space}

\bookgraphic{0.353}{part-2-14-0001-007-01}{}

\begin{verbatim}
  候选空间在观察评估结果前已经确定。
\end{verbatim}

\textbf{Adaptive Search Space}

\bookgraphic{0.353}{part-2-14-0001-007-02}{}

\begin{verbatim}
  新候选版本可能根据历史评估反馈自适应生成。
\end{verbatim}

这里必须区分两个概念：

\begin{readertable}
\begin{longtable}[]{@{}
  >{\raggedright\arraybackslash}p{(\columnwidth - 2\tabcolsep) * \real{0.37}}
  >{\raggedright\arraybackslash}p{(\columnwidth - 2\tabcolsep) * \real{0.59}}@{}}
\toprule
Quantity & Meaning \\
\midrule
\endhead
\(M=|\mathcal H|\) & 事先确定的完整假设空间大小 \\
\(K+1\) & RSI 实际访问的版本数量 \\
\bottomrule
\end{longtable}
\end{readertable}

\textbf{不能直接用 \(K+1\) 替代 \(M\)，然后声称得到了 RSI 的 Generalization Guarantee。}

因为实际生成的这些版本可能已经依赖于评估数据，从而破坏了逐个应用 Hoeffding 所需要的独立性条件。

不过，如果所有自适应生成的版本都属于一个预先确定的有限空间 \(\mathcal H\)，那么针对整个 \(\mathcal H\) 建立的 Uniform Bound 仍然有效。

这一区别是理解 Adaptive Overfitting 的关键。

\hypertarget{part-2-14-0001-ux4e00ux4e2aux53cdux4f8bux4e3aux4ec0ux4e48ux5019ux9009ux6a21ux578bux6570ux91cfux5c11ux4e5fux53efux80fdux4e25ux91cdux8fc7ux62dfux5408}{%
\subsubsection{8.2 一个反例：为什么候选模型数量少也可能严重过拟合？}\label{part-2-14-0001-ux4e00ux4e2aux53cdux4f8bux4e3aux4ec0ux4e48ux5019ux9009ux6a21ux578bux6570ux91cfux5c11ux4e5fux53efux80fdux4e25ux91cdux8fc7ux62dfux5408}}

设：
\[X\sim\operatorname{Uniform}(0,1),
\qquad Y=0.\]

观察到评估数据：
\[D_n=\{X_1,\ldots,X_n\}.\]

现在根据这个数据集构造一个分类器：
\[h_{D_n}(x)=
\begin{cases}
0,&x\in D_n,\\
1,&x\notin D_n.
\end{cases}\]

采用 Zero-One Loss。

显然，在评估数据集上：
\[\widehat R_n(h_{D_n})=0.\]

但是，由于连续型随机变量落在有限个固定点上的概率为零：
\[\mathbb P(X\in D_n\mid D_n)=0.\]

所以：
\[R(h_{D_n})=1.\]

最终：
\[\boxed{
G_n(h_{D_n})
=
R(h_{D_n})-\widehat R_n(h_{D_n})
=
1
}\]

即使我们最终只输出了一个模型，也可以产生最大的 Generalization Gap。

这个例子虽然极端，但严格说明了：

\textbf{实际输出的模型数量，并不足以刻画模型选择过程的统计复杂度。}

这也解释了为什么研究 RSI 时，不能只统计迭代轮数，还需要研究 Agent 从评估环境中获取了多少信息，以及如何利用这些信息。

\hypertarget{part-2-14-0001-ux91d1ux878dux573aux666fux4e2dux7684ux5176ux4ed6ux6761ux4ef6ux9650ux5236}{%
\subsubsection{8.3 金融场景中的其他条件限制}\label{part-2-14-0001-ux91d1ux878dux573aux666fux4e2dux7684ux5176ux4ed6ux6761ux4ef6ux9650ux5236}}

\begin{readertable}
\begin{longtable}[]{@{}
  >{\raggedright\arraybackslash}p{(\columnwidth - 2\tabcolsep) * \real{0.43}}
  >{\raggedright\arraybackslash}p{(\columnwidth - 2\tabcolsep) * \real{0.53}}@{}}
\toprule
Hoeffding Assumption & Potential Violation in Finance \\
\midrule
\endhead
Independent Samples & 金融时间序列具有时间依赖性 \\
Bounded Loss & Squared Error、Returns 可能无界 \\
Fixed Target Distribution & Non-stationarity、Regime Shift \\
Fixed Hypothesis Space & Agent 可能进行开放式、自适应搜索 \\
Sample Mean Estimator & Sharpe Ratio 是非线性统计量 \\
\bottomrule
\end{longtable}
\end{readertable}

因此，我们目前得到的是一个 \textbf{Theoretical Baseline}，而不是能够直接用于真实金融 RSI 系统的完整理论。

后续的 Adaptive Data Analysis 和 Time-Series Evaluation，正是为了处理这些条件受到破坏之后的问题。

\hypertarget{part-2-14-0001-research-connections-x-program-161}{%
\subsection{9. Research Connections --- X Program \#161}\label{part-2-14-0001-research-connections-x-program-161}}

经过本阶段的学习，我们可以进一步提出几个具体的实验方向。

\begin{readertable}
\begin{longtable}[]{@{}
  >{\raggedright\arraybackslash}p{(\columnwidth - 4\tabcolsep) * \real{0.25}}
  >{\raggedright\arraybackslash}p{(\columnwidth - 4\tabcolsep) * \real{0.34}}
  >{\raggedright\arraybackslash}p{(\columnwidth - 4\tabcolsep) * \real{0.37}}@{}}
\toprule
Research Variable & Experimental Design & Hypothesis \\
\midrule
\endhead
Evaluation Sample Size & 改变评估集规模 \(n\) & 更大的评估集可能缓解选择偏差 \\
Search Budget & 改变候选版本数量 & 更大的搜索预算可能扩大 Gap \\
Evaluation Noise & 控制 Synthetic Data 的噪声强度 & 高噪声可能加剧错误选择 \\
Search Adaptivity & 比较 Fixed Search 与 Adaptive Search & 自适应反馈可能带来额外的过拟合风险 \\
Search Space & 限制 Agent 允许修改的组件 & 限制搜索复杂度可能改善泛化 \\
\bottomrule
\end{longtable}
\end{readertable}

其中一个非常自然的实验就是：

保持底层模型、数据生成机制和评估指标不变，只改变 Sample Size 与 Search Budget。

随后观察：
\[G_k=S_k^{\mathrm{IS}}-S_k^{\mathrm{OOS}}\]

如何随着 RSI Iteration 变化。

我们可以利用前面的理论作为解释框架，但不能预先假设实验一定满足：
\[G_k\propto
\sqrt{\frac{\log M}{n}}.\]

因为真实的自适应搜索不一定满足该公式的成立条件。

\hypertarget{part-2-14-0001-ii-b-conclusion-index}{%
\subsection{10. II-B --- Conclusion Index}\label{part-2-14-0001-ii-b-conclusion-index}}

按照我们之前确定的学习方式，以下是本阶段需要保留的结论。推导可以通过 Index 返回相应章节。

\begin{readertable}
\begin{longtable}[]{@{}
  >{\raggedright\arraybackslash}p{(\columnwidth - 4\tabcolsep) * \real{0.11}}
  >{\raggedright\arraybackslash}p{(\columnwidth - 4\tabcolsep) * \real{0.71}}
  >{\raggedright\arraybackslash}p{(\columnwidth - 4\tabcolsep) * \real{0.14}}@{}}
\toprule
ID & Key Result & Index \\
\midrule
\endhead
B1 & LLN 提供渐近保证，Concentration Inequalities 提供有限样本保证 & §2 \\
B2 & Hoeffding's Inequality 提供指数型概率上界 & §3 \\
B3 & 固定模型的 Generalization Bound & §4 \\
B4 & Hoeffding + Union Bound 得到 Uniform Generalization Bound & §5.2--3 \\
B5 & 泛化误差与样本量、假设空间大小之间的关系 & §5.4 \\
B6 & Sample Complexity 随 \(\epsilon^{-2}\) 增长 & §6 \\
B7 & ERM Excess Risk 的有限样本保证 & §7 \\
B8 & 实际访问的候选数量不能直接替代固定假设空间大小 & §8.1 \\
B9 & Data-dependent Model 可以在候选数量极少时严重过拟合 & §8.2 \\
B10 & 经典泛化理论不能未经修正直接应用于金融 RSI & §8.3 \\
\bottomrule
\end{longtable}
\end{readertable}

\hypertarget{part-2-14-0001-core-formula-sheet}{%
\subsubsection{10.1 Core Formula Sheet}\label{part-2-14-0001-core-formula-sheet}}

本阶段最重要的是以下五个公式。

\textbf{Result 1 --- Single-model Generalization Bound}

\[\boxed{
\mathbb P(|\widehat R_n(h)-R(h)|>\epsilon)
\leq2e^{-2n\epsilon^2}
}\]

\textbf{Result 2 --- Uniform Generalization Bound}

\[\boxed{
\mathbb P
\left(
\sup_{h\in\mathcal H}
|\widehat R_n(h)-R(h)|>\epsilon
\right)
\leq2Me^{-2n\epsilon^2}
}\]

\textbf{Result 3 --- High-probability Error Bound}

\[\boxed{
\sup_{h\in\mathcal H}
|\widehat R_n(h)-R(h)|
\leq
\sqrt{\frac{\log(2M/\alpha)}{2n}}
}\]

\textbf{Result 4 --- Sample Complexity}

\[\boxed{
n\geq
\frac{\log(2M/\alpha)}{2\epsilon^2}
}\]

\textbf{Result 5 --- Approximate ERM Excess Risk Bound}

\[\boxed{
R(\tilde h_n)-R(h^*)
\leq
\sqrt{\frac{2\log(2M/\alpha)}n}+\delta
}\]

这些公式均有明确的成立条件。对于金融时间序列、无界损失及自适应生成的假设空间，需要额外的理论处理。

\hypertarget{part-2-14-0001-knowledge-checklist}{%
\subsubsection{10.2 Knowledge Checklist}\label{part-2-14-0001-knowledge-checklist}}

\textbf{II-B --- Knowledge Review}

\begin{itemize}
\item
  理解 LLN 与 Concentration Inequalities 的区别
\item
  理解 Hoeffding's Inequality 的成立条件
\item
  能够利用 Union Bound 推导 Uniform Generalization Bound
\item
  理解 Sample Size 与 Hypothesis Space Size 的数学关系
\item
  理解 Sample Complexity 的含义
\item
  理解 Statistical Error 与 Optimization Error 的联系
\item
  理解为什么实际 RSI 迭代次数不能直接代入 Uniform Bound
\item
  理解经典泛化理论在金融 RSI 中的局限性
\end{itemize}

\hypertarget{part-2-14-0001-next-ii-c-model-selection-bias-and-multiple-testing}{%
\subsection{Next --- II-C: Model Selection Bias and Multiple Testing}\label{part-2-14-0001-next-ii-c-model-selection-bias-and-multiple-testing}}

II-A --- Risk and Generalization

II-B --- Concentration Inequalities

II-C --- Model Selection Bias \& Multiple Testing

II-D --- Adaptive Overfitting \& Holdout Reuse

II-E --- Data Leakage \& Distribution Shift

II-F --- Time-Series Validation \& Backtesting

前面的 II-A 和 II-B 主要解决了评估误差的定义与控制问题。

下一阶段 II-C 将正式研究 \textbf{Selection Mechanism 本身如何产生系统性偏差}，包括 Winner's Curse、Order Statistics、Multiple Hypothesis Testing、FWER、FDR，以及金融领域中尤其重要的 Data Snooping。

我们将逐渐从研究泛化误差的上界，转向研究实际过拟合是如何产生的。
